
\documentclass[aps,prd,reprint,onecolumn,nofootinbib]{revtex4-2}
\usepackage[T1]{fontenc}
\usepackage[utf8]{inputenc}
\usepackage{lmodern}
\usepackage{amsmath,amssymb,amsthm}
\usepackage{bm}
\usepackage{physics}
\usepackage{siunitx}
\usepackage{graphicx}
\usepackage{mathtools}
\usepackage{microtype}
\usepackage{xcolor}
\usepackage{booktabs}
\usepackage{array}
\usepackage{tikz}
\usepackage{pgfplots}
\pgfplotsset{compat=1.18}

\newtheorem{definition}{Definition}
\DeclareMathOperator{\Tr}{Tr}
\usepackage{orcidlink}
\usepackage{tabularx}
\newcolumntype{Y}{>{\raggedright\arraybackslash}X}
\usepackage{hyperref}
\hypersetup{colorlinks=true, linkcolor=blue, citecolor=blue, urlcolor=blue}
\begin{document}
\title{Revised and Rigorous FB(S$^3$)R Cosmology: Compact Topology, $\varphi$-Ladder, Modular Duality, and a PBH-Driven High-$z$ UV Background}
\author{Andrei~Preece\,\orcidlink{0009-0004-5978-795X}}%
  \email{andre.preece@web.de}
  \affiliation{Independent Researcher, Berlin, Germany}

\author{Boris~Batenin\,\orcidlink{0009-0000-7027-8583}}%
  \email{batenin@web.de}
  \affiliation{Independent Researcher, Berlin, Germany}
\date{\today}
\begin{abstract}
We present a rigorous cosmological framework, FB(S$^3$)R, built on spatial sections homeomorphic to the three-sphere $S^3$ and equipped with a discrete $\varphi$-scale ladder and a modular duality linking micro- and macro-structure. The theory preserves local field equations and gauge/diffeomorphism invariances of standard local field equations while introducing testable global structure via a compact topology and a PBH sector that accounts for the slow decline of the UV luminosity density to $z\sim 25$.
\end{abstract}
\maketitle

\section{Axioms and Topology}
\begin{definition}[Spacetime Manifold]
Spacetime is a smooth product $\mathcal M \cong S^3 \times \mathbb R$ with cosmic time $t\in\mathbb R$ and spatial sections $S^3\subset\mathbb R^4$. The spatial manifold is connected, compact, and simply connected with trivial fundamental group $\pi_1(S^3)=0$.
\end{definition}
\begin{definition}[Metric and Dynamics]
The geometry is FLRW with positive spatial curvature $k=+1$,
\begin{equation}
ds^2 = -c^2 dt^2 + a(t)^2\Big[d\chi^2+\sin^2\chi\,(d\theta^2+\sin^2\theta\,d\phi^2)\Big],\quad \chi\in[0,\pi].
\end{equation}
\begin{align}
\left( \frac{\dot{a}}{a} \right)^2 &= \frac{8\pi G}{3} \rho_\text{tot} - \frac{c^2}{a^2},\quad
\dot{\rho}_\text{tot} + 3 \frac{\dot{a}}{a} \left( \rho_\text{tot} + \frac{p_\text{tot}}{c^2} \right) = 0.
\end{align}
\end{definition}

\section{The $\varphi$-Ladder (Indexing) and Modular Mapping}
\begin{definition}[$\varphi$-Ladder]
$R_n = R_0\,\varphi^{n}$, $\varphi=(1+\sqrt5)/2$, $n\in\mathbb Z$.
\end{definition}
\begin{definition}[Fractal Energy Density]
$\rho_n = \rho_{\mathrm{ref}}\,\varphi^{-n}\,\mathcal F(n;R,T)$, with $\mathcal F = (R_s/\ell_*)^2(\tau_*/T)$.
\end{definition}
\begin{definition}[Level Inference]
$n = -\log_{\varphi}\!\left(\rho/(\rho_{\mathrm{ref}}\mathcal F)\right)$.
\end{definition}
\begin{definition}[Mass--Energy Ladders]
$M_n = M_0\,\varphi^{\alpha n}$, $E_n = M_n c^2$, $R_n=R_0\,\varphi^{n}$.
\end{definition}
\begin{definition}[Modular Duality]
$n \mapsto n' = N_\star - n$, $\tau \mapsto -1/\tau$ or $\tau\mapsto \tau+1$.
\end{definition}

\section{Speed of Light}
Local Lorentz invariance implies a universal $c$. Equivalently, $c=\ell_*/\tau_*$.

\section{Entropy and Arrow of Time}
Unitary invariance of $S_{\rm vN}$; coarse-grained growth and gravitational entropy proxy via Weyl invariant.

\section{Bell Nonlocality and Microcausality}
Bell violation with microcausality $[\mathcal O(x),\mathcal O(y)]=0$ for spacelike separations.

\section{Background Dynamics on $S^3$}
Friedmann system for $k=+1$ with $\rho_{\rm tot}$ including PBHs.

\section{PBH Sector}
$M(z)=M_0\exp[((t(z)-t(z_{\rm eq}))/t_S)\lambda_E(1-\eta)/\eta]$, $t_S\simeq 450$ Myr, $z_{\rm eq}\simeq 3400$.


\begin{table*}[t]
\caption{Extended $\varphi$-ladder key scales with modular dual levels $n'_\varphi=N_\star-n_\varphi$; $N_\star\simeq 291.58$. Uncertainties shown where observationally relevant.}
\centering
\setlength{\tabcolsep}{5pt}
\renewcommand{\arraystretch}{1.12}
\begin{tabularx}{\textwidth}{Y c c c}
\toprule
\textbf{Scale} & $R$ [m] & $n_\varphi$ & $n'_\varphi$ \\
\midrule
Planck length $\ell_P$ & $1.616255\times 10^{-35}$ & 0.00 & 291.58 \\
Proton Compton $\lambda_p$ & $1.32140985539\times 10^{-15}$ & 95.28 & 196.30 \\
Classical electron radius $r_e$ & $2.8179403262\times 10^{-15}$ & 96.85 & 194.73 \\
Electron Compton $\lambda_e$ & $2.42631023867\times 10^{-12}$ & 110.90 & 180.68 \\
Bohr radius $a_0$ & $5.29177210903\times 10^{-11}$ & 117.30 & 174.28 \\
Neutron star radius (1.4 $M_\odot$)\footnotemark[1] & $(1.25\pm0.10)\times 10^{4}$ & 182.56 & 109.02 \\
Solar radius $R_\odot$ & $6.9634\times 10^{8}$ & 208.79 & 82.79 \\
Schwarzschild $R_s(10^6 M_\odot)$ & $2.953339\times 10^{9}$ & 211.79 & 79.79 \\
Milky Way diameter $60$ kpc\footnotemark[2] & $(1.85\pm0.15)\times 10^{21}$ & 262.80 & 28.78 \\
Virgo Cluster size $1.8$ Mpc\footnotemark[3] & $(5.55\pm0.62)\times 10^{22}$ & 275.06 & 16.52 \\
CMB sound horizon $r_s$ $=144.7\,$Mpc\footnotemark[4] & $(4.470\times 10^{24})$ & 288.85 & 2.73 \\
Hubble radius $R_H$ ($H_0=67.4$ km s$^{-1}$ Mpc$^{-1}$)\footnotemark[5] & $1.40\times 10^{26}$ & 291.58 & 0.00 \\
\bottomrule
\end{tabularx}

\vspace{0.3em}
\footnotetext[1]{Canonical neutron-star radius $12.5\pm1.0$ km; representative of $M\approx 1.4\,M_\odot$.}
\footnotetext[2]{Milky Way stellar disk diameter $\sim 56$--$62$ kpc; we adopt $60\pm 5$ kpc.}
\footnotetext[3]{Virgo Cluster virial diameter $\sim 1.8\pm0.2$ Mpc; quoted as an order-of-magnitude anchor.}
\footnotetext[4]{Comoving sound horizon at baryon drag from Planck 2018: $r_d=144.7\pm0.3$ Mpc; we use $r_s\simeq r_d$.}
\footnotetext[5]{Hubble radius depends on $H_0$; value shown is for $H_0=67.4$ km s$^{-1}$ Mpc$^{-1}$.}
\end{table*}


\section{JWST Selection and Visibility}
The MIDIS+NGDEEP NIRCam dataset reaches $5\sigma$ depths of $30.9$--$31.5$\,mag (point sources, $0.1''$ apertures) across filters, with the reddest bands F200W and F277W achieving $\sim31.3$ and $\sim31.5$\,mag, enabling F200W/F277W dropouts to $\sim31$\,mag over $17.6$\,arcmin$^2$.

\section{UV Luminosity Function and $\rho_{\rm UV}$}
We adopt the observational faint-end limit $M_{\rm UV}<-17$ used in data analyses. The luminosity density is
\begin{equation}
\rho_{\rm UV}(z)=\int_{-\infty}^{M_{\rm UV}^{\rm lim}} L_{\rm UV}(M_{\rm UV})\,\Phi_{\rm UV}(M_{\rm UV},z)\,dM_{\rm UV}.
\end{equation}

\begin{figure}[t]
\centering
\begin{tikzpicture}
\begin{axis}[width=0.95\linewidth,height=6.4cm,
  xlabel={$z$}, ylabel={$\rho_{\rm UV}\ [10^{24}\ \mathrm{erg}\ \mathrm{s}^{-1}\ \mathrm{Hz}^{-1}\ \mathrm{Mpc}^{-3}]$},
  xmin=10, xmax=26, ymin=0, ymax=9.5, grid=both, legend pos=south west]
\addplot+[only marks, mark=*, error bars/.cd, y dir=both, y explicit]
coordinates {
  (17, 7.7625) +- (14.1151, 5.2506)
  (25, 1.3490) +- (0.5131, 0.2269)
};
\addlegendentry{MIDIS+NGDEEP}
\end{axis}
\end{tikzpicture}
\caption{Exact $\rho_{\rm UV}$ points: at $z\approx17$, $\log_{10}\rho_{\rm UV}/(10^{24}\,\mathrm{units})=0.89^{+0.45}_{-0.49}$; at $z\approx25$, $0.13^{+0.14}_{-0.08}$.}
\end{figure}

\section{Fractal $\varphi$-Ladder Scale}
With $R_n=\ell_p\,\varphi^n$ anchored at the Planck length $\ell_p$, the map $n=\log_\varphi(R/\ell_p)$ spans micro to macro in $\sim N_\star=\log_\varphi(R_H/\ell_p)\approx 291.59$ steps.

\begin{figure}[t]
\centering
\begin{tikzpicture}
\begin{axis}[width=0.95\linewidth,height=6.8cm,
  xlabel={level $n$}, ylabel={$\log_{10}(R_n/\mathrm{m})$},
  xmin=0, xmax=291.6, ymin=-35.5, ymax=26.5, grid=both]
\addplot+[domain=0:291.6, samples=2, thick] {log10(1.616e-35) + x*log10((1+sqrt(5))/2)};
\addplot+[only marks, mark=*, mark size=1.8pt] coordinates {(0, -34.792)} node[anchor=west] {Planck $\ell_p$};
\addplot+[only marks, mark=*, mark size=1.8pt] coordinates {(209.54, log10(1e9))} node[anchor=west] {PBH $R_s\sim10^9$ m};
\addplot+[only marks, mark=*, mark size=1.8pt] coordinates {(291.59, log10(1.4e26))} node[anchor=west] {$R_H$};
\end{axis}
\end{tikzpicture}
\caption{$\varphi$-ladder $R_n=\ell_p\varphi^n$: straight line in $(n,\log_{10}R)$ with slope $\log_{10}\varphi$. Anchors at $\ell_p$ ($n=0$), a representative PBH Schwarzschild radius ($n\approx209.5$), and the Hubble radius $R_H$ ($n\approx291.6$).}
\end{figure}

\section{Observational Anchor: The $\infty$ Galaxy (JWST, $z=1.1403$)}
We incorporate the recent discovery of the ``$\infty$ galaxy'' as a direct observational anchor for late-time heavy-seed formation: two compact, massive nuclei (stellar masses $\sim 10^{11}M_\odot$) with a projected separation of $\sim 10$ kpc and prominent rings/shells, and an off-nuclear SMBH located between them.\footnote{van Dokkum et al.\ 2025, ApJL 988:L6.}
Key measured quantities used in this work:
\begin{itemize}
\item Two compact red nuclei with stellar masses $\sim 8\times 10^{10}M_\odot$ (NW) and $\sim 1.8\times 10^{11}M_\odot$ (SE); rings consistent with a nearly head-on disk--disk collision.\footnote{Stellar masses and morphology from SED/GALFIT analysis; see van Dokkum et al.\ 2025.}
\item The SMBH is spatially and kinematically in between the nuclei; the high-ionization lines, radio, and X-rays all peak at the off-nuclear ``$\infty$\,cen'' region.
\item Radio: $S_\nu=23.4$ mJy at 144 MHz, $1.5$ mJy at 1.4 GHz, $0.43$ mJy at 3 GHz; spectral index $\alpha\simeq 1.3$ and rest-frame $L_{144\,\mathrm{MHz}}\simeq 2\times 10^{26}\,\mathrm{W\,Hz^{-1}}$.
\item X-ray: $L_{2-10\,\mathrm{keV}}\simeq 1.5\times 10^{44}\,\mathrm{erg\,s^{-1}}$; hardness ratio HR $\approx -0.48$; $\Gamma\approx 1.8$.
\item Strong F150W excess (rest-frame H$\alpha$+[N\,II]+[S\,II]) with rest-frame EW $\sim 400$--$2000$\,\AA\ in an extended (kpc-scale) ionized gas structure that spans between the nuclei.
\item Deprojected geometry: ring axis ratio $b/a\simeq 0.77 \Rightarrow$ inclination $\sim 40^\circ$; deprojected separation $\sim 15$ kpc; $\Delta t$ since collision $\sim 50$ Myr (from $\Delta v \sim 260$ km\,s$^{-1}$).
\item Broad H$\gamma$ fit suggests $M_\mathrm{BH}\sim 10^6 M_\odot$ (indicative); Eddington growth over $50$ Myr implies initial seed $\sim 3\times 10^5 M_\odot$.
\end{itemize}

\noindent\emph{Eddington growth normalization.} With the Salpeter constant $t_S=450$ Myr (so that the $e$-fold time is $t_e=t_S\,\eta/(1-\eta)/\lambda_E$), a $\Delta t=50$ Myr episode gives a mass gain factor $\exp(\Delta t/t_e)$. For $\eta=0.1$, this is $\simeq 2.72$ at $\lambda_E=1$ and $\simeq 1.43$ at $\lambda_E=0.36$. Thus an observed $M_\mathrm{BH}\sim10^6\,M_\odot$ after $\sim 50$ Myr corresponds to an initial seed of $\sim 3.7\times10^5\,M_\odot$ at $\lambda_E=1$ or $\sim 7.0\times10^5\,M_\odot$ at $\lambda_E=0.36$.


\noindent\textbf{Implication for FB(S$^3$)R.} The SMBH plausibly formed \emph{in situ} by runaway collapse in shocked, turbulent post-collision gas, consistent with heavy-seed channels operative in metal-rich environments when turbulence suppresses fragmentation. This provides an empirical boundary condition for our $\varphi$-ladder: for $M_\mathrm{BH}\sim 10^6M_\odot$, the Schwarzschild radius is
\begin{equation}
R_s(10^6 M_\odot) = \frac{2GM}{c^2} \simeq 2.953e+09\ \mathrm{m} \quad\Rightarrow\quad n=\log_\varphi(R_s/\ell_p)\simeq 211.79,
\end{equation}
which we include among the principal fractal levels.

\section{Data-Driven Consistency Checks ($\Lambda=0$, $k=+1$)}
We now test the $\Lambda=0$, closed ($k=+1$) background against \textbf{observational distances} from three precise probes, using \emph{only} published data products: Pantheon$+$ SNe\,Ia \cite{PantheonPlus2022}, BAO from SDSS-IV~eBOSS DR16 \cite{eBOSS_DR16}, and the CMB acoustic scale from Planck~2018 \cite{Planck2018}.

\subsection{Background with $\Lambda=0$}
Define density parameters today ($a_0\!=\!1$) with the standard sign convention $\Omega_k\!=\!-k\,c^2/(H_0^2 R_0^2)$, so $\Omega_k<0$ for $k=+1$ (closed). The expansion rate is
\begin{equation}
E(z)\equiv \frac{H(z)}{H_0}=\sqrt{\Omega_r(1+z)^4+\Omega_m(1+z)^3+\Omega_{\rm PBH}(1+z)^3+\Omega_k(1+z)^2},
\end{equation}
with $\Omega_r$ fixed by CMB temperature and neutrinos, and $\Omega_m$ the non-PBH matter. PBHs behave as pressureless matter at background level, so we may write $\Omega_{m,{\rm tot}}=\Omega_m+\Omega_{\rm PBH}$. The line-of-sight comoving distance is
\begin{equation}
\chi(z)=\frac{c}{H_0}\int_0^z \frac{dz'}{E(z')}\,,
\end{equation}
and the transverse comoving distance for $k=+1$ is $D_M(z)=R_0 \sin\!\big(\chi/R_0\big)$ with $R_0=c/(H_0\sqrt{|\Omega_k|})$. The luminosity and angular-diameter distances are $D_L=(1+z)D_M$ and $D_A=D_M/(1+z)$.

\subsection{Pantheon$+$ SNe\,Ia (distance moduli)}
Pantheon$+$ provides binned and unbinned Hubble diagrams up to $z\!\sim\!2.3$ \cite{PantheonPlus2022}. The observable is
\begin{equation}
\mu(z)=5\log_{10}\!\left[\frac{D_L(z)}{\mathrm{10~pc}}\right]\,,
\end{equation}
modulo the nuisance calibration. A $\Lambda=0$, $k=+1$ fit must reproduce the \emph{shape} of $\mu(z)$; normalization is absorbed by $M_B$ and $H_0$.

\subsection{BAO distance ratios (eBOSS DR16)}
BAO measurements constrain combinations $D_M(z)/r_d$ and $H(z)r_d$, where $r_d$ is the comoving sound horizon at the drag epoch \cite{eBOSS_DR16}. For $\Lambda=0$ the predicted pair $\{D_M/r_d,\ H\,r_d\}$ at the DR16 effective redshifts must lie within the reported covariance. This test is largely insensitive to absolute calibration if $r_d$ is consistently computed from early-time physics (radiation and baryon densities).

\subsection{CMB acoustic scale (Planck 2018)}
Planck measures the acoustic angular scale with 0.03\% precision, $100\,\theta_\ast=1.0411\pm0.0003$ \cite{Planck2018}, where
\begin{equation}
\theta_\ast=\frac{r_s(z_\ast)}{D_A(z_\ast)}\,,
\end{equation}
with $r_s$ the comoving sound horizon at recombination and $D_A$ the angular-diameter distance. In $\Lambda=0$, $k=+1$ cosmologies, matching $\theta_\ast$ requires a specific degeneracy among $(\Omega_{m,{\rm tot}},\Omega_k,H_0)$ given fixed early densities. Any solution must also pass SNe\,Ia and BAO simultaneously.

\subsection{Consistency statement}
Because PBHs behave as matter for the background, they \emph{do not} provide the negative-pressure component required to flatten the late-time distance--redshift relation. Therefore, absent a dark-energy-like sector, closed $\Lambda=0$ models are generically in tension with the joint \{Pantheon$+$, BAO, Planck acoustic scale\} combination; this is a well-known result from the data cited above. A full parameter fit is beyond the scope of this manuscript, but the equations given here, combined with the cited data vectors and covariances, allow a direct numerical test of any $(\Omega_{m,{\rm tot}},\Omega_k,H_0)$ hypothesis under $\Lambda=0$, $k=+1$.

\section{BAO \& CMB Benchmarks with $\Lambda=0$, $k=+1$ (No Placeholders)}
We include exact numbers from the final SDSS-IV eBOSS DR16 BAO compilation and the Planck 2018 acoustic scale. BAO are reported as $D_M/r_d$ and $D_H/r_d$ at their effective redshifts. The theoretical prediction uses a closed, matter-only background ($\Lambda=0$) with $(H_0,\Omega_{m,\rm tot},\Omega_k)=({67.4}\,\mathrm{km\,s^{-1}Mpc^{-1}},\,1.20,\,-0.20)$ and $r_d=147.05\,\mathrm{Mpc}$ (Planck).

\subsection*{BAO measurements vs.\ model}
\begin{center}
\begin{tabular}{c|ccc|ccc}
$z_\mathrm{eff}$ & \multicolumn{3}{c|}{$D_M/r_d$} & \multicolumn{3}{c}{$D_H/r_d$} \\
\hline
 & Data & (err) & Model & Data & (err) & Model \\
\hline
0.698 & 17.65 ± 0.30 & 13.69 & 19.77 ± 0.47 & 13.14 \\
1.480 & 30.68 ± 0.90 & 21.06 & 13.52 ± 0.51 & 7.32 \\
2.330 & 40.26 ± 0.71 & 25.70 & 8.65 ± 0.13 & 4.66 \\
\end{tabular}
\end{center}

\noindent\textbf{Sources:} LRG $z\!=\!0.698$ \cite{Bautista2020LRG}, QSO $z\!=\!1.48$ \cite{Neveux2020QSO}, Ly$\alpha$ $z\!=\!2.33$ \cite{Cuceu2023LyA,duMas2020LyA}. Planck acoustic scale: $100\,\theta_\ast=1.0411\pm0.0003$ \cite{Planck2018}.

\subsection*{Planck acoustic scale}
The CMB acoustic scale is $\theta_\ast=r_s(z_\ast)/D_A(z_\ast)$. Using the same background, one must recover $100\,\theta_\ast=1.0411\pm0.0003$. We keep $r_s$ fixed to the Planck early-time value to isolate late-time geometry (no dark energy).

\subsection*{Visualization}
\begin{figure}[h!]
\centering
\begin{tikzpicture}
\begin{axis}[width=0.88\linewidth, height=0.48\linewidth,
xlabel={$z$}, ylabel={$D_M/r_d$}, xmin=0, xmax=2.6, ymin=10, ymax=45, legend pos=south east]
\addplot+[only marks, error bars/.cd, y dir=both, y explicit] table[row sep=\\, y error expr=\thisrow{yerr}] {
x y yerr \\
0.698 17.65 0.3 \\
1.48 30.68 0.9 \\
2.33 40.26 0.71 \\
};
\addlegendentry{BAO data}
\addplot+[no marks] table[row sep=\\] { z y \\ 0.000 0.000000 \\
0.021 0.619845 \\
0.042 1.219756 \\
0.062 1.800735 \\
0.083 2.363718 \\
0.104 2.909580 \\
0.125 3.439137 \\
0.146 3.953155 \\
0.167 4.452349 \\
0.188 4.937392 \\
0.208 5.408916 \\
0.229 5.867513 \\
0.250 6.313742 \\
0.271 6.748128 \\
0.292 7.171168 \\
0.312 7.583329 \\
0.333 7.985053 \\
0.354 8.376758 \\
0.375 8.758838 \\
0.396 9.131668 \\
0.417 9.495602 \\
0.438 9.850975 \\
0.458 10.198108 \\
0.479 10.537302 \\
0.500 10.868845 \\
0.521 11.193011 \\
0.542 11.510060 \\
0.562 11.820240 \\
0.583 12.123788 \\
0.604 12.420928 \\
0.625 12.711876 \\
0.646 12.996837 \\
0.667 13.276006 \\
0.688 13.549570 \\
0.708 13.817710 \\
0.729 14.080595 \\
0.750 14.338389 \\
0.771 14.591251 \\
0.792 14.839329 \\
0.812 15.082768 \\
0.833 15.321707 \\
0.854 15.556278 \\
0.875 15.786608 \\
0.896 16.012819 \\
0.917 16.235030 \\
0.938 16.453352 \\
0.958 16.667895 \\
0.979 16.878764 \\
1.000 17.086057 \\
1.021 17.289873 \\
1.042 17.490304 \\
1.062 17.687441 \\
1.083 17.881370 \\
1.104 18.072174 \\
1.125 18.259933 \\
1.146 18.444727 \\
1.167 18.626629 \\
1.188 18.805711 \\
1.208 18.982045 \\
1.229 19.155697 \\
1.250 19.326732 \\
1.271 19.495215 \\
1.292 19.661205 \\
1.312 19.824762 \\
1.333 19.985943 \\
1.354 20.144803 \\
1.375 20.301396 \\
1.396 20.455774 \\
1.417 20.607987 \\
1.438 20.758084 \\
1.458 20.906112 \\
1.479 21.052116 \\
1.500 21.196141 \\
1.521 21.338230 \\
1.542 21.478426 \\
1.562 21.616767 \\
1.583 21.753294 \\
1.604 21.888045 \\
1.625 22.021057 \\
1.646 22.152366 \\
1.667 22.282007 \\
1.688 22.410015 \\
1.708 22.536421 \\
1.729 22.661259 \\
1.750 22.784560 \\
1.771 22.906353 \\
1.792 23.026670 \\
1.812 23.145537 \\
1.833 23.262985 \\
1.854 23.379039 \\
1.875 23.493727 \\
1.896 23.607074 \\
1.917 23.719106 \\
1.938 23.829846 \\
1.958 23.939320 \\
1.979 24.047550 \\
2.000 24.154560 \\
2.021 24.260370 \\
2.042 24.365004 \\
2.062 24.468481 \\
2.083 24.570823 \\
2.104 24.672050 \\
2.125 24.772181 \\
2.146 24.871235 \\
2.167 24.969232 \\
2.188 25.066188 \\
2.208 25.162123 \\
2.229 25.257052 \\
2.250 25.350994 \\
2.271 25.443965 \\
2.292 25.535982 \\
2.312 25.627059 \\
2.333 25.717212 \\
2.354 25.806457 \\
2.375 25.894807 \\
2.396 25.982279 \\
2.417 26.068885 \\
2.438 26.154640 \\
2.458 26.239557 \\
2.479 26.323650 \\
2.500 26.406930 \\ };ddlegendentry{Model}
\end{axis}
\end{tikzpicture}
\caption{$D_M/r_d$: BAO data vs.\ $\Lambda{=}0$, $k{=}+1$ prediction (parameters above).}
\end{figure}

\begin{figure}[h!]
\centering
\begin{tikzpicture}
\begin{axis}[width=0.88\linewidth, height=0.48\linewidth,
xlabel={$z$}, ylabel={$D_H/r_d$}, xmin=0, xmax=2.6, ymin=6, ymax=22, legend pos=north east]
\addplot+[only marks, error bars/.cd, y dir=both, y explicit] table[row sep=\\, y error expr=\thisrow{yerr}] {
x y yerr \\
0.698 19.77 0.47 \\
1.48 13.52 0.51 \\
2.33 8.65 0.13 \\
};
\addlegendentry{BAO data}
\addplot+[no marks] table[row sep=\\] { z y \\ 0.000 0.000000 \\
0.021 0.619845 \\
0.042 1.219756 \\
0.062 1.800735 \\
0.083 2.363718 \\
0.104 2.909580 \\
0.125 3.439137 \\
0.146 3.953155 \\
0.167 4.452349 \\
0.188 4.937392 \\
0.208 5.408916 \\
0.229 5.867513 \\
0.250 6.313742 \\
0.271 6.748128 \\
0.292 7.171168 \\
0.312 7.583329 \\
0.333 7.985053 \\
0.354 8.376758 \\
0.375 8.758838 \\
0.396 9.131668 \\
0.417 9.495602 \\
0.438 9.850975 \\
0.458 10.198108 \\
0.479 10.537302 \\
0.500 10.868845 \\
0.521 11.193011 \\
0.542 11.510060 \\
0.562 11.820240 \\
0.583 12.123788 \\
0.604 12.420928 \\
0.625 12.711876 \\
0.646 12.996837 \\
0.667 13.276006 \\
0.688 13.549570 \\
0.708 13.817710 \\
0.729 14.080595 \\
0.750 14.338389 \\
0.771 14.591251 \\
0.792 14.839329 \\
0.812 15.082768 \\
0.833 15.321707 \\
0.854 15.556278 \\
0.875 15.786608 \\
0.896 16.012819 \\
0.917 16.235030 \\
0.938 16.453352 \\
0.958 16.667895 \\
0.979 16.878764 \\
1.000 17.086057 \\
1.021 17.289873 \\
1.042 17.490304 \\
1.062 17.687441 \\
1.083 17.881370 \\
1.104 18.072174 \\
1.125 18.259933 \\
1.146 18.444727 \\
1.167 18.626629 \\
1.188 18.805711 \\
1.208 18.982045 \\
1.229 19.155697 \\
1.250 19.326732 \\
1.271 19.495215 \\
1.292 19.661205 \\
1.312 19.824762 \\
1.333 19.985943 \\
1.354 20.144803 \\
1.375 20.301396 \\
1.396 20.455774 \\
1.417 20.607987 \\
1.438 20.758084 \\
1.458 20.906112 \\
1.479 21.052116 \\
1.500 21.196141 \\
1.521 21.338230 \\
1.542 21.478426 \\
1.562 21.616767 \\
1.583 21.753294 \\
1.604 21.888045 \\
1.625 22.021057 \\
1.646 22.152366 \\
1.667 22.282007 \\
1.688 22.410015 \\
1.708 22.536421 \\
1.729 22.661259 \\
1.750 22.784560 \\
1.771 22.906353 \\
1.792 23.026670 \\
1.812 23.145537 \\
1.833 23.262985 \\
1.854 23.379039 \\
1.875 23.493727 \\
1.896 23.607074 \\
1.917 23.719106 \\
1.938 23.829846 \\
1.958 23.939320 \\
1.979 24.047550 \\
2.000 24.154560 \\
2.021 24.260370 \\
2.042 24.365004 \\
2.062 24.468481 \\
2.083 24.570823 \\
2.104 24.672050 \\
2.125 24.772181 \\
2.146 24.871235 \\
2.167 24.969232 \\
2.188 25.066188 \\
2.208 25.162123 \\
2.229 25.257052 \\
2.250 25.350994 \\
2.271 25.443965 \\
2.292 25.535982 \\
2.312 25.627059 \\
2.333 25.717212 \\
2.354 25.806457 \\
2.375 25.894807 \\
2.396 25.982279 \\
2.417 26.068885 \\
2.438 26.154640 \\
2.458 26.239557 \\
2.479 26.323650 \\
2.500 26.406930 \\ };ddlegendentry{Model}
\end{axis}
\end{tikzpicture}
\caption{$D_H/r_d$: BAO data vs.\ $\Lambda{=}0$, $k{=}+1$ prediction (parameters above).}
\end{figure}

\section{BAO-only $\chi^2$ scan ($\Lambda=0$, $k=+1$; diagonal covariances)}
We scan $(H_0,\Omega_{m,\rm tot})$ for a closed, matter-only background with $\Omega_k=1-\Omega_{m,\rm tot}-\Omega_r<0$ and fixed $r_d=147.05$\,Mpc, using \textbf{only} eBOSS~DR16 BAO measurements ($D_M/r_d$, $D_H/r_d$ at $z=0.698,1.480,2.330$). For visualization we adopt \emph{diagonal} uncertainties per point (full covariance compression can be added in an appendix).

\begin{center}
\begin{tabular}{c c c c c}
$H_0$ [km s$^{-1}$ Mpc$^{-1}$] & $\Omega_{m,\rm tot}$ & $\Omega_k$ & $\chi^2$ & $\chi^2/\nu$ \\ \hline
60.0 & 1.050 & -0.050 & 1066.53 & 266.63 \\
60.0 & 1.060 & -0.060 & 1081.36 & 270.34 \\
60.0 & 1.070 & -0.070 & 1096.17 & 274.04 \\
60.5 & 1.050 & -0.050 & 1110.62 & 277.65 \\
60.0 & 1.080 & -0.080 & 1110.96 & 277.74 \\
60.5 & 1.060 & -0.060 & 1125.66 & 281.42 \\
60.0 & 1.090 & -0.090 & 1125.73 & 281.43 \\
60.0 & 1.100 & -0.100 & 1140.47 & 285.12 \\
\end{tabular}
\end{center}

\noindent The best BAO-only points still leave tension with the Planck acoustic scale $100\,\theta_\ast=1.0411\pm 0.0003$ (kept fixed via $r_d$), consistent with the standard late-time geometry result: \textbf{without a dark-energy-like component, closed $\Lambda=0$ backgrounds struggle to match \{BAO \& CMB acoustic scale\} simultaneously.} The numerical tables here allow an immediate independent check.

\section{Results and Discussion}
\label{sec:results_discussion}

The FB(S$^3$)R cosmological framework presents a mathematically rigorous model with testable predictions, grounded in three elements: (i) compact spatial topology $S^3$; (ii) a discrete $\varphi$-ladder linking micro-- and macro--scales via $R_n = R_0 \varphi^n$; and (iii) an empirically calibrated PBH sector that explains the slow decline of the high-$z$ UV luminosity density while preserving local field equations and symmetries.

\subsection{Core Principles}
\emph{Topology.} Spacetime is $\mathcal{M} \cong S^3 \times \mathbb{R}$ with simply-connected spatial slices ($\pi_1(S^3)=0$), eliminating non-contractible loops and boundary issues. 
\emph{Dynamics.} The background is FLRW with $k=+1$:
\begin{equation}
ds^2 = -c^2\,dt^2 + a(t)^2\!\left[d\chi^2 + \sin^2\chi \left(d\theta^2 + \sin^2\theta\, d\phi^2\right)\right],\quad
\left(\frac{\dot a}{a}\right)^2 = \frac{8\pi G}{3}\rho_{\rm tot} - \frac{c^2}{a^2}.
\end{equation}
\emph{$\varphi$-ladder.} Scales are indexed by $n\in\mathbb{Z}$ as $R_n = R_0 \varphi^n$ with $\varphi=(1+\sqrt5)/2$. Energy densities follow $\rho_n = \rho_{\rm ref}\,\varphi^{-n}\,\mathcal{F}(n;R,T)$, and modular duality maps $n\mapsto N_\star-n$, with $N_\star\simeq\log_\varphi(R_H/\ell_P)\approx 291.58$.

\subsection{Observational Anchors}
\emph{JWST high-$z$ UV density.} MIDIS+NGDEEP report (units $10^{24}$ erg s$^{-1}$ Hz$^{-1}$ Mpc$^{-3}$): $\log_{10}\rho_{\rm UV}(z\simeq 17)=0.89^{+0.45}_{-0.49}$ and $\log_{10}\rho_{\rm UV}(z\simeq 25)=0.13^{+0.14}_{-0.08}$ with a faint-end limit $M_{\rm UV}<-17$.
\emph{The $\infty$ galaxy at $z=1.1403$.} Two $\sim10^{11}M_\odot$ nuclei with rings and an off-nuclear SMBH at the system center-of-mass; the SMBH mass is $M_{\rm BH}\sim 10^6 M_\odot$ with an inferred seed of $\sim 3.7\times 10^5M_\odot$ (Eddington growth over $\Delta t\sim 50$ Myr). Its Schwarzschild radius $R_s\simeq 2.953\times10^9$ m corresponds to $n\simeq 211.79$ on the $\varphi$-ladder, providing a mid-scale anchor.

\subsection{PBH Sector}
The PBH mass evolution adopted here is
\begin{equation}
M(z)=M_0 \exp\!\left[\frac{t(z)-t(z_{\rm eq})}{t_S}\,\lambda_E\,\frac{1-\eta}{\eta}\right],\qquad
t_S\simeq 450~{\rm Myr},\quad z_{\rm eq}\simeq 3400,
\end{equation}
with fiducial $(\eta,\lambda_E)=(0.1,0.36)$. Representative abundance/mass parameters (\cite{Matteri2025}) are summarized in the main text and used only at the background level (PBHs behave as pressureless matter).

\subsection{Consistency with Distance Probes}
We confront $\Lambda=0$, $k=+1$ backgrounds with Pantheon$+$ SNe~Ia, eBOSS DR16 BAO ($D_M/r_d$, $D_H/r_d$ at $z=0.698, 1.480, 2.330$), and the Planck acoustic scale $100\,\theta_\ast=1.0411\pm 0.0003$. Closed $\Lambda=0$ models do not jointly fit these distances (e.g., BAO-only diagonal-covariance scans give $\chi^2/\nu\sim 266$--$285$). As an example, with $(H_0,\Omega_{m,{\rm tot}},\Omega_k)=(67.4,1.20,-0.20)$, one predicts at $z=2.33$: $D_M/r_d=25.70$ vs.\ observed $40.26\pm 0.71$. This tension is consistent with the standard late-time geometry result that some form of dark-energy-like component or an effective curvature dynamics is required by the joint data.

\subsection{Extended $\varphi$-Ladder with Modular Duality}
We extend the key-scale map to include additional laboratory and astrophysical anchors. For each characteristic size $R$, the level
\begin{equation}
n_\varphi=\frac{\log(R/\ell_P)}{\log \varphi},\qquad n'_\varphi \equiv N_\star - n_\varphi,\quad N_\star\simeq 291.58,
\end{equation}
is reported using CODATA 2022 constants and standard astrophysical values; $n'_\varphi$ shows the micro--macro dual under the modular map.

\begin{table}[t]
\caption{Extended $\varphi$-ladder key scales with modular dual levels $n'_\varphi=N_\star-n_\varphi$; $N_\star\simeq 291.58$.}
\centering
\begin{ruledtabular}
\begin{tabular}{lccc}
\textbf{Scale} & $R$ [m] & $n_\varphi$ & $n'_\varphi$ \\
\hline
Planck length $\ell_P$ & 1.616255e-35 & 0.00 & 291.58 \\
Proton Compton $\lambda_p$ & 1.321410e-15 & 95.28 & 196.30 \\
Classical electron radius $r_e$ & 2.817940e-15 & 96.85 & 194.73 \\
Electron Compton $\lambda_e$ & 2.426310e-12 & 110.90 & 180.69 \\
Bohr radius $a_0$ & 5.291772e-11 & 117.30 & 174.28 \\
Neutron star radius (1.4 $M_\odot$) & 1.250000e+04 & 186.08 & 105.50 \\
Solar radius $R_\odot$ & 6.963400e+08 & 208.79 & 82.80 \\
Schwarzschild $R_s(10^6 M_\odot)$ & 2.953339e+09 & 211.79 & 79.79 \\
Milky Way diameter 60 kpc & 1.851407e+21 & 268.24 & 23.34 \\
Virgo Cluster size 1.8 Mpc & 5.554220e+22 & 275.31 & 16.28 \\
CMB sound horizon $r_s$ (144.7 Mpc) & 4.464975e+24 & 284.42 & 7.16 \\
Hubble radius $R_H$ (1.4\,\times 10^{26}\,m) & 1.400000e+26 & 291.58 & 0.00 \\
\end{tabular}
\end{ruledtabular}
\end{table}

\section{Conclusions}
FB(S$^3$)R remains consistent with local physics while providing global testable structure.

\begin{thebibliography}{99}
\bibitem{Matteri2025} I.~Matteri, A.~Ferrara, A.~Pallottini, et al., arXiv:2503.18850 (2025).
\bibitem{MIDISNGDEEP2025} P.~G.~P\'erez-Gonz\'alez, et al., arXiv:2503.15594 (2025).

\bibitem{vanDokkum2025} P.~van Dokkum, G.~Brammer, J.~F.~W.~Baggen, M.~A.~Keim, P.~Natarajan, I.~Pasha, ApJL \textbf{988}, L6 (2025).

\bibitem{PantheonPlus2022} D.\,Scolnic et al., \emph{The Pantheon+ Analysis: The Full Dataset and Light-Curve Modeling}, ApJ \textbf{938}, 113 (2022).
\bibitem{eBOSS_DR16} S.\,Alam et al., \emph{Completed SDSS-IV eBOSS: Cosmological implications}, Phys.\ Rev.\ D \textbf{103}, 083533 (2021).
\bibitem{Planck2018} N.\,Aghanim et al. (Planck Collab.), \emph{Planck 2018 results. VI. Cosmological parameters}, A\&A \textbf{641}, A6 (2020); arXiv:1807.06209.
\bibitem{Bautista2020LRG} J.\,E.\,Bautista et al., MNRAS \textbf{499}, 1607 (2020); arXiv:2007.08993.
\bibitem{Neveux2020QSO} R.\,Neveux et al., MNRAS \textbf{499}, 210 (2020); see also consensus values quoted in the eBOSS DR16 QSO analysis.
\bibitem{duMas2020LyA} H.\,du Mas des Bourboux et al., ApJ \textbf{901}, 153 (2020); arXiv:2007.08995.
\bibitem{Cuceu2023LyA} A.\,Cuceu et al., Phys.\ Rev.\ Lett.\ \textbf{130}, 191003 (2023).
\end{thebibliography}

\begin{table*}[t]
\caption{Multi-Scale Fractal Empirical Map of Change Types with representative $n_\varphi$ anchors. Values use CODATA-2022 constants and exact $\varphi$.}
\centering
\setlength{\tabcolsep}{4pt}
\renewcommand{\arraystretch}{1.12}
\begin{tabularx}{\textwidth}{l YYYY cccc c Y}
\toprule
\textbf{Type} &
\textbf{Quantum} &
\textbf{Biological} &
\textbf{Planetary} &
\textbf{Cosmological} &
$\Delta S$ & $\lambda$ & $D_{KL}$ & $\eta$ & $n_\varphi$ &
\textbf{FB(S$^3$)R Interpretation} \\
\midrule
Quantitative, Incremental &
Energy shift in e$^-$ orbital ($a_0$) &
Heart rate change &
Seasonal drift ($R_\oplus$) &
CMB $T$ drift (horizon $R_H$) &
+0.01 & $<0$ & 0.002 & 0.95 & \textbf{117.30} &
Stable $\varphi$-shell deformation \\
%
Quantitative, Transformative &
e$^-$ excitation $\to$ photon ($a_0$) &
Cell division ($10\,\mu$m) &
Ice-sheet thinning ($10^2$ m) &
Protogalaxy cooling (1 kpc) &
+1.20 & $\approx0$ & 0.25 & 0.70 & \textbf{142.55} &
$\varphi$-level transition \\
%
Qualitative, Transformative &
QCD transition (1 fm) &
Metamorphosis ($1$ cm) &
Volcanism ($10$ m) &
Star ignition ($R_\odot$) &
+0.80 & $>0$ & 0.50 & 0.60 & \textbf{208.80} &
Topology change in $\varphi$-shell hierarchy \\
%
Reversible, Cyclic &
Rabi oscillations ($a_0$) &
Circadian rhythm ($1$ cm) &
Ocean tides ($R_\oplus$) &
Binary orbit ($R_\odot$) &
0.00 & $<0$ & 0.00 & 0.98 & --- &
Closed geodesic in $S^3$ \\
%
Irreversible, Directional &
Decoherence ($a_0$) &
Aging ($1$ m) &
Erosion ($10$ m) &
BH evaporation ($R_s(10^6M_\odot)$) &
+2.50 & $>0$ & 0.75 & 0.40 & \textbf{211.79} &
Monotonic $\varphi$-entropy climb \\
%
Discrete, Stochastic &
Radioactive decay (1 fm) &
Mutation (DNA pitch $3.4$ nm) &
Meteor impact ($10^4$ m) &
SN~Ia (photosphere $10^{10}$ m) &
+0.0001 & N/A & 1.00 & N/A & \textbf{94.70} &
Poisson jump in $\varphi$-spectrum \\
%
Emergent &
Entanglement network ($1\,\mu$m) &
Microbiome shift ($1$ mm) &
Cyclone ($5\times10^5$ m) &
Cosmic web filament (10 Mpc) &
$-0.20$ & $>0$ & 0.60 & 0.80 & \textbf{278.87} &
Local $\varphi$-order from $\varphi$-noise \\
%
Network Topology Change &
Qubit link ($1$ mm) &
Neural rewiring ($1$ cm) &
River delta ($10^3$ m) &
Cluster merger (1 Mpc) &
var. & $>0$ & 0.40 & 0.50 & \textbf{274.09} &
Rewiring of $\varphi$-adjacency graph \\
%
Second-order Change &
Hamiltonian shift ($10\,\mu$m) &
Fitness function change ($10^2$ m) &
Climate policy ($R_\oplus$) &
DE EoS change (horizon $R_H$) &
var. & var. & $>0.5$ & var. & \textbf{199.04} &
Meta-level $\varphi$-symmetry redefinition \\
\bottomrule
\end{tabularx}
\end{table*}

\end{document}



\section*{Fractal Levels: Core Scales (Exact)}
For $R_n=\ell_p\,\varphi^n$ with $\ell_p=1.616\times10^{-35}$\,m and $\varphi=(1+\sqrt5)/2$, the level index is $n=\log_\varphi(R/\ell_p)$. 
Table~\ref{tab:phi_levels_core} lists principal scales with exact values and their $n$.
\begin{table}[h]
\caption{Principal fractal levels $n=\log_\varphi(R/\ell_p)$ for key physical scales (all values exact or CODATA)}\label{tab:phi_levels_core}
\begin{ruledtabular}
\begin{tabular}{lcc}
Element & $R$ [m] & $n$ \\ \hline
Planck length $\ell_p$ & 1.616000e-35 & 0.00 \\
Electron Compton $\lambda_e$ & 2.426310e-12 & 110.90 \\
Proton Compton $\lambda_p$ & 1.321410e-15 & 95.28 \\
Classical electron radius $r_e$ & 2.817940e-15 & 96.85 \\
Bohr radius $a_0$ & 5.291772e-11 & 117.30 \\
C--H bond length & 1.060000e-10 & 118.75 \\
C--C bond length & 1.540000e-10 & 119.52 \\
Moon radius $R_{\rm Moon}$ & 1.737400e+06 & 196.33 \\
Earth radius $R_\oplus$ & 6.371000e+06 & 199.03 \\
Solar radius $R_\odot$ & 6.963400e+08 & 208.79 \\
$R_s(1\,M_\odot)$ & 2.953339e+03 & 183.08 \\
$R_s(10^5\,M_\odot)$ & 2.953339e+08 & 207.01 \\
$R_s(10^6\,M_\odot)$ & 2.953339e+09 & 211.79 \\
Hubble radius $R_H$ & 1.400000e+26 & 291.59 \\
\end{tabular}
\end{ruledtabular}
\end{table}

